%% ch15 --- Życie, rynki i szum: jak odróżnić chaos od przypadku
%%
%% Napisany we wrześniu 2026. Każda liczba, której czytelnik nie policzy w
%% pamięci, pochodzi z code/measure/ch15.py, który importuje TE SAME funkcje,
%% które drukują listingi (code/ch15/). Liczby losowe pochodzą z generatora
%% Pythona z nazwanymi ziarnami, więc tasowania i szum są takie same na każdym
%% komputerze; błędy prognoz są drukowane z dwiema cyframi znaczącymi.
%%
%% Czego ten rozdział nie mówi: patrz nagłówek bliźniaka.
%%
%% Bliźniak: chapters/en/ch15-life-noise.tex. Podrozdział za podrozdziałem,
%% makro za makrem, liczba za liczbą; tools/parity.py je porównuje.

\chapter{Życie, rynki i szum: jak odróżnić chaos od przypadku}
\label{ch:life}
\index{chaos!odróżnianie od szumu}\index{szum}

\noindent
Kardiolog patrzy na odstępy między uderzeniami serca, epidemiolog na sto lat
comiesięcznych zachorowań na odrę, inwestor na cenę akcji. Każdy widzi
poszarpaną linię, która nigdy dokładnie się nie powtarza, i chce wiedzieć, co
nią szarpie: reguła, w której kilka liczb popycha się nawzajem, czy przypadek,
czyli niezliczone drobne przyczyny, których żaden model nigdy nie wymieni.
Najczęściej jedno i drugie.

Rozdział~\ref{ch:logistic} pokazał, że reguła długości jednej linijki potrafi
dawać wyniki nieregularne jak losowanie totolotka. Ten rozdział zadaje to
pytanie od drugiej strony: gdy ktoś poda ci poszarpany szereg liczb bez
żadnej reguły, czy umiesz stwierdzić, czy zrobiła go reguła? Zbudujesz dwa
testy, zmierzysz, jak dobrze działają na danych zaszumionych i krótkich, a
potem odwiedzisz miejsca, w których ogłaszano odkrycie chaosu. Przesłanie nie
brzmi: chaos jest wszędzie, ani: nie ma go nigdzie. Brzmi: twierdzenie o
chaosie da się sprawdzić.

\begin{outcomes}
\outcome{zrobić szereg chaotyczny i szereg losowy z dokładnie tych samych
  liczb i powiedzieć, dlaczego żaden histogram ich nie rozróżni}
\outcome{narysować odwzorowanie powrotu, policzyć kratki, w które trafia, i
  odczytać ten wynik}
\outcome{prognozować szereg z jego własnej przeszłości metodą analogów i
  zmierzyć, jak błąd rośnie z liczbą kroków naprzód}
\outcome{powiedzieć, jak szum pomiarowy i krótki zapis osłabiają oba testy}
\outcome{ocenić twierdzenie o chaosie w populacjach, epidemiach, sercach czy
  na rynkach na podstawie przedstawionych dowodów}
\end{outcomes}

\section{Dwa szeregi, jeden zestaw liczb}
\label{sec:ch15-twins}
\index{surogat}

\noindent
Zacznij od zagadki. Dwa szeregi zawierają po \val{ch15.n} liczb z przedziału
od $0$ do $1$. Jeden zrobiła reguła logistyczna z rozdziału~\ref{ch:logistic}
z pokrętłem ustawionym na $r = 4$: jeśli $x_n$ to wartość w kroku $n$, to
następna wynosi
\[ x_{n+1} = 4\,x_n\,(1 - x_n), \]
czyli cztery razy wartość razy jeden minus wartość. Drugi szereg jest losowy.

\begin{think}
Szereg~A zaczyna się od \val{ch15.a.1}, \val{ch15.a.2}, \val{ch15.a.3},
\val{ch15.a.4}. Szereg~B zaczyna się od \val{ch15.b.1}, \val{ch15.b.2},
\val{ch15.b.3}, \val{ch15.b.4}. Który zrobiła reguła? Czy da się to poznać
na oko, a czy da się to policzyć?
\end{think}
\reveal{B. Na oko się nie da. Rachunkiem tak: $4 \times \num{0.2} \times
\num{0.8} = \num{0.64}$, czyli druga wartość B, a $4 \times \num{0.64} \times
\num{0.36}$ to \val{ch15.b.3}, jej trzecia. Dla A reguła daje z pierwszej
wartości około \val{ch15.a.next}, a druga wartość wynosi \val{ch15.a.2}.}
Zwróć uwagę, czego użył ten test. Nie wielkości liczb ani tego, jak często
pojawia się każda wielkość, tylko ich \emph{kolejności}: czy każda wartość
wynika z poprzedniej. O tym jest cały ten rozdział.

\begin{trapbox}
\textbf{\enquote{Ciąg, który wygląda na losowy, musi pochodzić z losowości.}}
Szereg~B wyglądał na równie losowy jak A, a powstał z jednej linijki
arytmetyki bez grama przypadku. To, że coś wygląda losowo, jest faktem o
wyglądzie. Nie mówi nic o przyczynie, dopóki nie sprawdzisz, czy jest w nim
porządek, jaki zostawia po sobie reguła.
\end{trapbox}

\noindent
Szereg~A nie powstał z rzutów kostką. Powstał z szeregu~B przez
przetasowanie:

\pyregion{code/ch15/twins.py}{make}{Dwa szeregi z jednego: reguła i te same liczby przetasowane}{lst:ch15-make}

\noindent
Pierwsza funkcja iteruje regułę; druga tasuje kopię listy generatorem liczb
losowych Pythona, uruchomionym od ustalonego \emph{ziarna}, dzięki czemu
tasowanie jest takie samo na każdym komputerze, a twój wynik zgadza się ze
stroną. Szereg zbudowany z innego przez pomieszanie kolejności nazywa się
\emph{surogatem}: to losowy szereg najtrudniejszy do odróżnienia od
oryginału, bo ma dokładnie te same liczby.

\begin{think}
Jak histogram przetasowanego szeregu \dash{} ile jego wartości wpada do każdej
piątej części przedziału od $0$ do $1$ \dash{} wypadnie w porównaniu z
histogramem szeregu reguły?
\end{think}
\reveal{Identycznie, słupek w słupek. Histogram liczy wartości i pomija ich
kolejność, a tasowanie zmienia tylko kolejność.}

\transcript{ch15-twins}

\noindent
Spójrz na kształt: \val{ch15.hist.low} spośród \val{ch15.n} wartości leży w
najniższej piątej części, a tylko \val{ch15.hist.mid} w środkowej, bo reguła
lubi przebywać blisko $0$ i $1$. Przetasowany szereg ma ten sam przekrzywiony
histogram, a także tę samą średnią, ten sam rozrzut, tę samą największą i
najmniejszą wartość. Żadne podsumowanie, które pomija kolejność, nie odróżni
tych bliźniaków.

\plotfig{ch15-twins}{Szereg reguły i jego przetasowanie. U góry: pierwsze sto
kroków każdego, równie poszarpane. U dołu: ich histogramy, identyczne.}{dwa
szeregi, jeden histogram}

\section{Każda wartość względem następnej}
\label{sec:ch15-return}
\index{odwzorowanie powrotu}

\noindent
W zagadce mogłeś sprawdzać każdą wartość względem poprzedniej, bo znałeś
regułę. Przy prawdziwych danych nikt ci jej nie poda, więc niech dane same ją
narysują. Postaw punkt w $(x_n, x_{n+1})$ dla każdego kroku $n$: w poziomie
ta wartość, w pionie następna. Taki rysunek nazywa się \emph{odwzorowaniem
powrotu}.

\begin{think}
Gdzie leżą punkty odwzorowania powrotu szeregu reguły? A przetasowanego?
\end{think}
\reveal{Punkty reguły leżą na jednej krzywej, garbie $y = 4x(1 - x)$, bo
każda następna wartość \emph{jest} regułą zastosowaną do poprzedniej. Punkty
przetasowania rozchodzą się po całym kwadracie, bo po przetasowaniu następna
wartość nie ma nic wspólnego z obecną.}

\plotfig{ch15-return-maps}{Każda wartość względem następnej. Szereg reguły
rysuje własną regułę; przetasowany szereg wypełnia kwadrat. Cienkie linie
dzielą kwadrat na sto małych kratek.}{krzywa i chmura}

\noindent
Rysunek to ocena; liczba to pomiar. Podziel kwadrat na sto małych kratek i
policz, w ile z nich trafia odwzorowanie powrotu:

\pyregion{code/ch15/return_map.py}{pairs}{Pary sąsiadów i kratki, w które trafiają}{lst:ch15-pairs}

\transcript{ch15-return-map}

\noindent
Krzywa może przejść tylko przez ograniczoną liczbę kratek. \val{ch15.n}
wartości reguły trafia w \val{ch15.sq.rule}, a dłuższy zapis nie zwiększyłby
tej liczby, bo każdy nowy punkt ląduje na tej samej krzywej. Przetasowanie
wciąż znajduje nowe kratki, aż zajmie wszystkie \val{ch15.sq.shuf}.

\begin{lab}{l15_01_return}{Krzywa czy chmura?}
Napisz test odwzorowania powrotu samodzielnie: połącz każdą wartość w parę z
następną, policz kratki, w które trafiają pary, i zwróć \code{"curve"}, gdy
zajętych jest mniej niż połowa kratek, a \code{"cloud"} w przeciwnym razie.
Potem wypróbuj siatkę dwa na dwa na szeregu, który krąży między trzema
wartościami, i zobacz, jak zbyt gruba siatka wprowadza w błąd.
\end{lab}

\begin{worldbox}
Robert Shaw, fizyk z Uniwersytetu Kalifornijskiego w Santa Cruz, mierzył
czas między kroplami spadającymi z cieknącego kranu. Przy niektórych
szybkościach przepływu krople spadały w nieregularnych odstępach, a gdy
zaznaczył każdy odstęp względem następnego, punkty nie wypełniły chmury:
ułożyły się wzdłuż cienkiego zakrzywionego kształtu. Kuchenny kran stał się
jednym z najsłynniejszych pokazów chaosu odczytanego z jednej mierzonej
wielkości.
\end{worldbox}

\noindent
Shaw mierzył jedną rzecz, odstęp między kroplami, choć w wodzie dzieje się
znacznie więcej, niż opisze jedna liczba. To, że w ogóle zobaczył kształt,
nie jest szczęściem. W 1980 roku Norman Packard, James Crutchfield, Doyne
Farmer i sam Shaw pokazali, jak odbudować kształt ruchu układu z jednej
mierzonej wielkości, używając jej opóźnionych kopii \dash{} odczytu teraz,
chwilę temu i jeszcze chwilę wcześniej \dash{} w zastępstwie liczb, których
nikt nie zmierzył. W 1981 roku Floris Takens udowodnił, że to działa.
Odwzorowanie powrotu to przypadek najprostszy: jedna wielkość, jedno
opóźnienie.

\begin{rigourbox}
Twierdzenie Takensa słowami: dla gładkiego układu deterministycznego, którego
stan ma $d$ liczb, $2d + 1$ opóźnionych kopii niemal dowolnego pojedynczego
pomiaru daje wierną kopię ruchu, wygiętą i rozciągniętą, ale o tym samym
kształcie. Dowód, w jego pracy z 1981 roku \emph{Detecting strange
attractors in turbulence}, korzysta z topologii różniczkowej.
\end{rigourbox}

\section{Prognoza z przeszłości}
\label{sec:ch15-analogues}
\index{metoda analogów}\index{najbliższy sąsiad}

\noindent
Drugi test zadaje pytanie, które obchodzi użytkownika danych: czy da się to
przewidzieć? Bez reguły możesz przeszukać zapis w poszukiwaniu chwili
najbardziej podobnej do obecnej i założyć, że to, co wtedy nastąpiło, nastąpi
i teraz. To \emph{metoda analogów}. W 1969 roku Edward Lorenz wypróbował ją na
mapach pogody i nie znalazł dnia z przeszłości dość podobnego, by się
przydał: przy tak wielu liczbach w stanie atmosfera nie wróciła w pobliże
miejsc, w których już była. Dla szeregu pojedynczych liczb metoda działa:

\pyregion{code/ch15/analogues.py}{forecast}{Prognoza według najbliższego analogu z przeszłości}{lst:ch15-forecast}

\noindent
\code{forecast} znajduje w przeszłości wartość najbliższą \code{x}, jej
\emph{najbliższego sąsiada}, i zwraca to, co nastąpiło \code{h} kroków po
nim. \code{mean\_error} uczy się na pierwszym tysiącu wartości, prognozuje
każdą wartość po nich i uśrednia wielkości chybień.

\begin{think}
Dla szeregu reguły średnie chybienie o krok naprzód wynosi około
\val{ch15.fc.rule}, czyli jedną tysięczną. Strzał na ślepo \dash{} dowolna
wartość wzięta z histogramu \dash{} chybia średnio o około \val{ch15.fc.shuf}.
Rozdział~\ref{ch:lyapunov} zmierzył, że przy $r = 4$ mały błąd mniej więcej
podwaja się w każdym kroku. O ile kroków naprzód może sięgnąć prognoza, zanim
jej chybienia dojdą do połowy chybień strzału na ślepo?
\end{think}
\reveal{Około dziewięciu: jedna tysięczna podwojona dziewięć razy to mniej
więcej połowa, bo $2^{9} = 512$. Pomiar się zgadza: chybienia po raz pierwszy
przekraczają połowę chybień przetasowania w kroku \val{ch15.fc.half}.}

\transcript{ch15-analogues}

\plotfig{ch15-forecast}{Średnie chybienie prognozy metodą analogów w
zależności od liczby kroków naprzód, w skali logarytmicznej. Przetasowanie to
strzał na ślepo przy każdym horyzoncie; szereg reguły prognozuje się dobrze i
z każdym krokiem gorzej.}{błąd prognozy a liczba kroków naprzód}

\noindent
Dla przetasowanego szeregu prognoza chybia o około \val{ch15.fc.shuf} przy
każdym horyzoncie: jego przeszłość nie niesie żadnej informacji o następnej
wartości. Dla szeregu reguły chybienie o krok naprzód jest około
\val{ch15.fc.times} razy mniejsze i rośnie mniej więcej \val{ch15.fc.double}
razy na krok \dash{} to rozciąganie, które rozdział~\ref{ch:lyapunov} zmierzył
od środka, widziane z zewnątrz \dash{} aż po kilkunastu krokach też staje się
strzałem na ślepo.

Chaos i szum różnią się więc w obie strony. Szum jest równie nieprzewidywalny
przy każdym horyzoncie. Chaos jest przewidywalny na kilka kroków naprzód,
gdzie szum nie jest, i traci tę przewidywalność w stałym, mierzalnym tempie:
to znacznie mocniejszy odcisk palca chaosu niż jakakolwiek nieregularność.

\section{Prawdziwe dane to mieszanina}
\label{sec:ch15-noisy}
\index{szum!pomiarowy}

\noindent
Nic nie jest mierzone dokładnie. Waga, termometr czy urzędnik przepisujący
rejestr przesuwają każdą wartość trochę, losowo. Żeby zobaczyć, co to robi,
przesuń każdą wartość szeregu reguły o losową wielkość nie większą niż pewien
rozmiar:

\pyregion{code/ch15/noisy.py}{noise}{Szum pomiarowy: każda wartość przesunięta losowo}{lst:ch15-noise}

\begin{think}
Dodaj do szeregu reguły szum o rozmiarze do \val{ch15.noise.small}: jedną
setną, mniej niż grubość krzywej na wydrukowanym odwzorowaniu powrotu. Co się
stanie ze średnim chybieniem o krok naprzód, które wynosiło \val{ch15.fc.rule}?
\end{think}
\reveal{Wzrośnie do \val{ch15.noise.small.err}, wielokrotnie. Analog, który
znajduje prognoza, sam jest przesunięty nawet o jedną setną, podobnie jak
prognozowana wartość, i to zanim podwajanie reguły w ogóle się zacznie.}

\transcript{ch15-noisy}

\noindent
W górnej tabeli, dla całego zapisu, krzywa w miarę wzrostu szumu grubieje w
pas: przy szumie do \val{ch15.noise.big}, jednej piątej całego zakresu,
zajmuje \val{ch15.noise.big.sq} kratki. Mimo to przy \val{ch15.n} wartościach
prognoza wciąż jest lepsza od strzału na ślepo: \val{ch15.noise.big.err}
wobec \val{ch15.fc.shuf}.

\plotfig{ch15-noisy}{Odwzorowania powrotu szeregu reguły oglądane przez szum
pomiarowy trzech rozmiarów. Krzywa grubieje w pas, a przy największym
rozmiarze pas prawie wypełnia kwadrat.}{krzywa grubiejąca w szumie}

Dolna tabela to część okrutna. Tnie zapis na krótkie kawałki i liczy, na sto
kawałków, jak często każdy test uznaje kawałek reguły za mniej losowy od jego
własnego przetasowania. Przy czterdziestu zaszumionych wartościach test
prognozy ma rację \val{ch15.short.fc.forty} razy. Przy dziesięciu ma rację
\val{ch15.short.fc} razy, a test kratek tylko \val{ch15.short.sq} \dash{}
gorzej niż rzut monetą. Krótkie, zaszumione dane potrafią sprawić, że porządek
wygląda jak przypadek, a przypadek jak porządek.
A to najłatwiejszy chaos, jaki istnieje, z jedną liczbą na krok; układ z
wieloma liczbami w stanie potrzebuje znacznie dłuższego zapisu.

\begin{lab}{l15_02_forecast}{Zapytaj kilka analogów}
Gdy w zapisie jest szum, najbliższy analog może być bliski przypadkiem.
Napisz prognozę, która znajduje kilka najbliższych wartości z przeszłości i
uśrednia to, co po nich nastąpiło. Test sprawdza ją na małej bibliotece,
którą przeliczysz ręcznie, a potem sprawdza, że uśrednienie pięciu analogów
jest lepsze od zaufania jednemu na zaszumionym szeregu.
\end{lab}

\section{Populacje, epidemie, serca}
\label{sec:ch15-life}
\index{dynamika populacji}\index{epidemia}\index{model SIR}

\noindent
Praca Roberta Maya z 1976 roku o odwzorowaniu logistycznym była pisana dla
ekologów. Jej lekcja brzmiała: populacja, której liczebność skacze
nieregularnie z roku na rok, nie musi być popychana przez pogodę czy
szczęście; może to robić jej własna reguła rozmnażania.

\begin{think}
Wiele prawdziwych populacji owadów rzeczywiście rośnie i maleje nieregularnie
z roku na rok. Czy to czyni je chaotycznymi?
\end{think}
\reveal{Nie. Nieregularne liczby daje chaos, ale daje je też populacja w
równowadze szarpana przez pogodę, choroby i przypadek. Nieregularność nie
odróżni tych dwóch; mogą to zrobić tylko testy porządku.}

\begin{trapbox}
\textbf{\enquote{Nieregularne dane dowodzą chaosu.}} Każdy szereg w tym
rozdziale był nieregularny, a połowa z nich to przetasowania. Chaos to
twierdzenie, że nieregularność tworzy reguła z niewieloma liczbami, i wymaga
dowodów na tę regułę: struktury w odwzorowaniu powrotu, prognozy, która jest
dobra na początku i pogarsza się równomiernie, oraz zapisu dość długiego i
czystego, by jedno i drugie coś znaczyło.
\end{trapbox}

\noindent
Najmocniejsze dowody przyszły z laboratorium, gdzie regułę można zmieniać
celowo. W 1997 roku Robert Costantino ze współpracownikami opisali w
\emph{Science} trojszyki, drobne chrząszcze żyjące w mące, hodowane w
słojach, którym eksperymentatorzy ręcznie ustalali tempo śmierci dorosłych i
dołączania nowych dorosłych. Gdy te tempa się zmieniały, kolonie przechodziły
od stałej liczebności przez cykle do nieregularnych wahań, tam gdzie
przewidywał ich model. Dzika populacja liczona raz do roku przez kilka
dziesięcioleci daje mniej więcej tak mało wartości jak dolna tabela powyżej.

Drugi klasyczny przykład to epidemie. \emph{Model SIR} dzieli ludzi na
podatnych, którzy mogą się zarazić, zakaźnych, którzy chorują, i
ozdrowieńców. Nowe zakażenia każdego dnia są proporcjonalne do tego, ilu jest
podatnych i zakaźnych, którzy mogą się spotkać, chory zaraża średnio przez
\val{ch15.sir.days} dni, a narodziny dostarczają nowych podatnych. Weź
chorobę podobną do odry, w której jeden chory wśród ludzi bez odporności
zaraża około \val{ch15.sir.r0} innych, i długość życia \val{ch15.sir.life}
lat. Potem dodaj rok szkolny: dzieci stykają się częściej w czasie nauki niż
w wakacje, więc niech tempo kontaktów rośnie i maleje raz w roku o ułamek
nazwany \emph{sezonowością}.

\pyregion{code/ch15/epidemic.py}{rule}{Epidemia z rokiem szkolnym: reguła SIR krokami po ćwierć doby}{lst:ch15-sir}

\transcript{ch15-epidemic}

\noindent
Bez roku szkolnego model się ustala: \val{ch15.sir.steady} zachorowań na
tysiąc osób, co roku tyle samo. Przy sezonowości \val{ch15.sir.two} lata duże
i małe występują na przemian, jak odra w Anglii i Walii, zanim w 1968 roku
wprowadzono szczepienia. Przy \val{ch15.sir.four} wzór powtarza się co cztery
lata, a przy \val{ch15.sir.wild} nie powtarza się wcale, a pchnięcie o jedną
miliardową na starcie posyła lata w inną stronę. To droga do chaosu z
rozdziału~\ref{ch:bifurcations}, przebyta przez epidemię.

Tu też potrzebna jest uczciwość. W chaotycznym ustawieniu epidemie modelu
wypalają się tak dokładnie, że w najniższym punkcie zakaźna jest mniej niż
jedna osoba na $10^{\val{ch15.sir.trough}}$. Żadne miasto nie ma milionowej
części człowieka: prawdziwa choroba by wygasła. Czy odra kiedykolwiek
była chaotyczna, wciąż jest przedmiotem sporu. W 1990 roku George Sugihara i
Robert May zastosowali prognozę metodą analogów do comiesięcznych zachorowań
na odrę w Nowym Jorku i znaleźli podpis z tego rozdziału, dobrą prognozę,
która równomiernie pogarsza się z horyzontem, podczas gdy ospa wietrzna w tym
samym mieście wyglądała jak roczny cykl z szumem na wierzchu.

Z sercem jest jeszcze trudniej. W doświadczeniach laboratoryjnych opisanych w
1981 roku skupiska komórek serca pobudzane regularnymi impulsami elektrycznymi
przechodziły przez podwojenie okresu do nieregularnego bicia. Zdrowe serce też
nie bije jak metronom, ale niegdyś popularne twierdzenie, że jego zmienność to
chaos, opiera się na testach z tego rozdziału zastosowanych do krótkich,
zaszumionych zapisów i pozostaje sporne. Zmienność to nie to samo co chaos.

\section{Rynki i czego wymaga twierdzenie o chaosie}
\label{sec:ch15-evidence}
\index{rynki}

\noindent
Ceny to najpilniej obserwowane poszarpane linie na świecie i przez pewien czas
wielu miało nadzieję, że to chaos jest wzorem ukrytym pod nimi.

\begin{think}
Przypuśćmy, że ktoś udowodnił, że cenę akcji tworzy chaotyczna reguła z
zaledwie kilkoma liczbami w stanie. Czy mógłbyś się wzbogacić na jej
prognozowaniu?
\end{think}
\reveal{Niewiele. Chaos da się prognozować tylko na kilka kroków naprzód,
tylko tak dobrze, jak pozwalają pomiary, i wcale poza swoim horyzontem. A
każdy inwestor, który znalazłby tę samą regułę, grałby według niej i zmieniał
ceny, które ona opisuje.}

\begin{trapbox}
\textbf{\enquote{Chaos na rynkach oznacza, że da się je przewidzieć.}}
Chaotyczne ceny obiecywałyby krótki horyzont z twardą granicą za nim, a nie
szklaną kulę. W rzeczywistości ekonomiści, którzy szukali, znaleźli stopy
zwrotu trudne do przewidzenia oraz okresy spokoju i burzy skupione razem, ale
dowody na chaotyczną regułę z niewieloma liczbami w stanie są słabe i sporne.
Nieregularne, skokowe i trudne do przewidzenia to razem jeszcze nie chaotyczne.
\end{trapbox}

\noindent
Zbierz cały rozdział, a twierdzenie o chaosie w prawdziwych danych wymaga:

\begin{itemize}
\item \textbf{długiego zapisu}, długiego w porównaniu z liczbą wielkości,
  które mają znaczenie;
\item \textbf{małego szumu} albo testów, które go uwzględniają;
\item \textbf{uczciwego porównania}: szereg musi wygrać ze swoimi surogatami,
  także z takimi, które zachowują jego powolne wzrosty i spadki, bo szereg z
  pamięcią, jak dobowa temperatura, daje się prognozować na dzień naprzód bez
  żadnego chaosu;
\item \textbf{właściwego kształtu błędu prognozy}: małego o krok naprzód i
  rosnącego w stałym tempie;
\item a najlepiej \textbf{mechanizmu}, najlepiej z pokrętłem, które da się
  przekręcić, jak u chrząszczy.
\end{itemize}

\noindent
Niewiele twierdzeń o sercach, mózgach czy gospodarkach spełnia wszystkie
pięć. To teoria chaosu wykonująca swoją pracę: mówi ci, do jakich wniosków
poszarpana linia cię uprawnia, a do jakich nie.

\begin{summarybox}
\begin{itemize}
\item \textbf{Surogat} zrobiony przez tasowanie ma te same liczby, więc ten
  sam \textbf{histogram}. Tylko \textbf{kolejność} odróżni regułę od
  przypadku.
\item \textbf{Odwzorowanie powrotu} przedstawia każdą wartość względem
  następnej: reguła rysuje krzywą, szum wypełnia kwadrat. \textbf{Zanurzenie
  z opóźnieniem} (Takens, 1981) odbudowuje kształt układu z jednej mierzonej
  wielkości.
\item \textbf{Metoda analogów} prognozuje według najbliższej wartości z
  przeszłości. Chaos prognozuje się dobrze o krok naprzód i gorzej w stałym
  tempie; szum jest strzałem na ślepo przy każdym horyzoncie.
\item \textbf{Szum pomiarowy} i \textbf{krótkie zapisy} osłabiają oba testy;
  przy dziesięciu zaszumionych wartościach test kratek wypada gorzej niż
  moneta.
\item Chaos przekonująco pokazano w populacjach laboratoryjnych; dla odry,
  serc i rynków dowody są mieszane, sporne albo słabe. \textbf{Nieregularne to
  nie chaotyczne}.
\item Twierdzenie o chaosie wymaga długiego, czystego zapisu, uczciwego
  porównania z surogatami, właściwego kształtu błędu prognozy, a najlepiej
  mechanizmu.
\end{itemize}
\end{summarybox}

\canyou

\begin{checkpoint}
\item Co tasowanie szeregu zachowuje, a co niszczy?
  \answerto{Zachowuje każdą wartość, więc i histogram. Niszczy kolejność, a to
  w niej widać by było regułę łączącą każdą wartość z następną.}
\item Odwzorowanie powrotu pewnego szeregu leży na jednej cienkiej krzywej. Co
  to mówi o tym, jak powstał szereg?
  \answerto{Że każdą wartość wyznacza poprzednia: zrobiła go reguła jednej
  liczby. Szum rozrzuciłby punkty po kwadracie.}
\item Jeden szereg zajmuje $30$ ze stu kratek przy $500$ wartościach i nadal
  $30$ przy $5000$. Inny zajmuje $80$, a potem $100$. Który zachowuje się jak
  szereg reguły?
  \answerto{Pierwszy: krzywa przechodzi tylko przez ograniczoną liczbę kratek,
  ile by punktów na niej nie wylądowało. Drugi wciąż wypełnia kwadrat, jak
  szum.}
\item Prognoza metodą analogów chybia o około \num{0.001} o krok naprzód, a
  jej chybienia podwajają się w każdym kroku. Po ilu mniej więcej krokach
  chybienie dojdzie do \num{0.25}?
  \answerto{Po około ośmiu: $2^{8} = 256$, a $\num{0.001} \times 256$ to
  około \num{0.26}.}
\item Dlaczego zapis z dziesięciu zaszumionych wartości jest słabym dowodem za
  chaosem albo przeciw niemu?
  \answerto{Oba testy są przy tak małej liczbie zawodne: w pomiarze z tego
  rozdziału test kratek na dziesięciu zaszumionych wartościach miał rację
  rzadziej niż rzucona moneta.}
\item Co dzieje się z modelową epidemią, gdy rok szkolny staje się silniejszy,
  i co czyni jej chaotyczne ustawienie wątpliwym obrazem prawdziwej odry?
  \answerto{Jej lata przechodzą od jednakowych, przez naprzemienne i
  powtarzające się co cztery lata, do niepowtarzających się wcale; a jej
  chaotyczne dołki są tak głębokie, że prawdziwa choroba by wygasła.}
\item Ktoś pokazuje ci, że cena akcji jest nieregularna. Czy to dowód chaosu?
  Co byłoby dowodem?
  \answerto{Nie: szum też jest nieregularny. Dowodem byłaby prognoza, która
  wygrywa z surogatami, z błędem, który zaczyna od małego i rośnie
  równomiernie, w długim, czystym zapisie.}
\end{checkpoint}
